% V2第20章原书思考题20.1--20.10，PDF物理页204，独立候选；图20.20--20.21待精确重绘。
\section*{第20章\quad 电磁感应：思考题解析}
\paragraph{思考题 20.1}
线圈运动使其在磁场中穿过的磁通改变，按 \(\mathcal E=-d\Phi_B/dt\) 诱导电流；短接回路电阻变小，感生电流较强。由楞次定律，电流产生磁力矩阻碍线圈运动，机械能转成 \(I^2R\) 焦耳热，形成 \(\boxed{\text{电磁阻尼}}\)。若开路，感应电流很小，磁阻尼减弱。

\paragraph{思考题 20.2}
交流线圈产生时变磁通，金属内形成涡电流 \(\boldsymbol J=\sigma\boldsymbol E\)；金属耗散功率密度 \(\boxed{p_J=\boldsymbol J\cdot\boldsymbol E=J^2/\sigma}\)。若为铁磁金属还会有磁滞损耗。能源来自驱动线圈的高频电源，不是“磁场凭空产生热”。高频趋肤效应改变发热分布。

\paragraph{思考题 20.3}
时变磁通在铁芯横截面形成涡电流回路，整体厚铁芯回路面积大、回路电阻低，涡流发热较强。采用相互绝缘的薄硅钢片，片层截断大面积横向闭合电流回路，使涡流损耗显著降低。叠片应使工作磁通主要\emph{沿片面}穿行，以保持所需磁路并阻止垂直片面的环流。铁磁材料的磁滞损耗另由材料选择控制。

\paragraph{思考题 20.4}
同一面积、同一时变磁通下法拉第电动势 \(\boxed{\mathcal E_{\rm Cu}=\mathcal E_{\rm Al}=-d\Phi_B/dt}\)，与线圈材质无关。电阻 \(R=\rho \ell/A_{\rm wire}\) 一般不同，准稳态且可忽略自感时 \(\boxed{I_{\rm Cu}/I_{\rm Al}=R_{\rm Al}/R_{\rm Cu}}\)。快速交流时回路满足 \(L\,dI/dt+RI=\mathcal E(t)\)，还须考虑相位、交流趋肤与几何。

\paragraph{思考题 20.5}
变化磁场在环形轨道上生成\emph{非保守}环向电场，满足 \(\oint\boldsymbol E\cdot d\boldsymbol l=-d\Phi_B/dt\)。电子沿切向获得电功 \(dK/dt=q\boldsymbol E\cdot\boldsymbol v\)，其动能来自外部磁场电源提供的电磁能；静磁洛伦兹力 \(q\boldsymbol v\times\boldsymbol B\) 本身不做功。

\paragraph{思考题 20.6}
两线圈的互感本质为一个线圈电流在另一个线圈中产生的通量与电流之比 \(M=\Phi_{21}/I_1\)。让中间线圈的面积法向尽量与两侧线圈通量方向垂直，使穿过其回路的净磁通近零；对细小近似均匀场的线圈，这意味着中间线圈平面接近与场方向平行。实际近场不均匀时零互感需要对称性和实验调整，不能仅靠轴心相同保证。

\paragraph{思考题 20.7}
互感与\emph{相互穿链磁通}有关。两环共轴且小环套入大环时，常可让大环磁力线较多地穿过小环；与面对面远距离放置相比通常更大。但严格最大互感取决于半径比、轴向间距、相对方向和绕组形状，不能只用“套入”作普遍结论；同向法向使 \(M\) 正，翻转法向改变符号。

\paragraph{思考题 20.8}
断电瞬间电流企图快速下降，电感产生 \(\boxed{\mathcal E_L=-L\,dI/dt}\)，力图维持原电流。断开速度越快，\(|dI/dt|\) 越大，触点间可能形成很高瞬态电压并击穿空气，储存的 \(\tfrac12 LI^2\) 磁能在火花与电路中释放。火花不是证明开关外部突然出现新的能源。

\paragraph{思考题 20.9}
理想超导态排斥内部磁感应（Meissner效应），界面屏蔽电流造成有效抗磁响应，在非均匀磁场中可产生向上的磁力。对II类超导体，磁通进入涡旋并被缺陷\emph{钉扎}，可稳定悬浮的位置和方向。单靠静态完美抗磁排斥与稳定悬浮并非完全等同，具体取决于几何和钉扎；图20.20需按实物装置核验。

\paragraph{思考题 20.10}
探头交流磁场穿入金属，使金属中产生涡电流，其二次磁场反过来改变探头线圈的等效阻抗、相位和感应信号。频率、导电率、磁导率及金属形状决定响应。若只有\emph{时间不变}直流磁场，静止金属内无持续的涡电流，无法用相同感应原理测量；但移动磁铁或金属时仍可产生时变通量。
